
\subsection{Der Spin als Folge der Theorie}
\label{subsec:spin_als_folge_der_theorie}

Der Spin\index{Spin!als Folge der Dirac-Gleichung} des Elektrons war
bereits vor der Dirac-Gleichung aus Experimenten bekannt. In der
nichtrelativistischen Quantenmechanik musste er jedoch als zusätzliche
Eigenschaft eingeführt werden.

Die Dirac-Gleichung änderte diese Situation grundlegend. Der Spin
erscheint in ihr nicht als nachträgliche Ergänzung, sondern ergibt sich
aus der mathematischen Struktur der relativistischen Theorie.

\begin{HistoryBox}[Vom experimentellen Befund zur Theorie]
	\label{box:spin_vom_experiment_zur_theorie}
	
	Experimente hatten gezeigt, dass das Elektron zwei mögliche
	Spinrichtungen besitzt und sich in einem Magnetfeld entsprechend
	aufspaltet.
	
	Die Schrödingergleichung allein konnte diese Beobachtung nicht
	erklären.
	
	Erst die Dirac-Gleichung zeigte, dass ein relativistisches Elektron
	zwangsläufig eine innere zweifache Struktur besitzt, die dem Spin
	\(\hbar/2\) entspricht.
\end{HistoryBox}

\subsubsection*{Drehimpuls in der klassischen Mechanik}

In der klassischen Mechanik entsteht ein Drehimpuls%
\index{Drehimpuls!klassischer} beispielsweise dann, wenn sich ein
Teilchen auf einer Kreisbahn bewegt.

Der Bahndrehimpuls lautet

\begin{equation}
	\mathbf{L}
	=
	\mathbf{r}
	\times
	\mathbf{p},
	\label{eq:bahndrehimpuls_spin_dirac}
\end{equation}

wobei

\begin{itemize}
	\item \(\mathbf{r}\) der Ortsvektor und
	\item \(\mathbf{p}\) der Impuls
\end{itemize}

ist.

Dieser Drehimpuls hängt von der Bewegung des Teilchens im Raum ab. Ein
Elektron in einem Atom kann deshalb einen Bahndrehimpuls besitzen.

Der Spin ist jedoch etwas anderes. Er ist unabhängig davon vorhanden,
ob sich das Elektron auf einer Bahn bewegt oder nicht.

\begin{DidacticBox}[Bahndrehimpuls und Spin]
	\label{box:bahndrehimpuls_und_spin}
	
	Der Bahndrehimpuls entsteht durch die Bewegung eines Teilchens im
	Raum:
	
	\begin{equation}
		\mathbf{L}
		=
		\mathbf{r}\times\mathbf{p}.
	\end{equation}
	
	Der Spin ist dagegen eine innere Quanteneigenschaft des Elektrons.
	
	Auch ein Elektron ohne Bahndrehimpuls besitzt Spin.
\end{DidacticBox}

\subsubsection*{Warum der Spin keine klassische Rotation ist}

Der Begriff Spin stammt vom englischen Wort für Drehung. Diese
Bezeichnung kann jedoch zu einem falschen Bild führen.

Das Elektron darf nicht als kleine geladene Kugel verstanden werden,
die sich um ihre eigene Achse dreht. Ein solches klassisches Modell
würde zu erheblichen Schwierigkeiten führen.

Wollte man das gemessene magnetische Moment durch eine rotierende
Kugel erklären, müssten sich Teile ihrer Oberfläche unter bestimmten
Annahmen schneller als das Licht bewegen. Außerdem gibt es keinen
experimentellen Hinweis auf eine räumliche Ausdehnung oder Oberfläche
des Elektrons.

Der Spin ist daher keine mechanische Eigenrotation, sondern eine
fundamentale Eigenschaft des Quantenzustandes.

\begin{DidacticBox}[Kein rotierendes Kügelchen]
	\label{box:spin_keine_klassische_rotation}
	
	Der Elektronenspin darf nicht als klassische Drehung eines kleinen
	Körpers verstanden werden.
	
	Der Begriff beschreibt eine quantenmechanische Form des Drehimpulses,
	die keine unmittelbare Entsprechung in der klassischen Mechanik
	besitzt.
\end{DidacticBox}

\subsubsection*{Die vierkomponentige Struktur}

Die Wellenfunktion der Dirac-Gleichung ist ein vierkomponentiger
Spinor\index{Dirac-Spinor!Spinstruktur}:

\begin{equation}
	\Psi
	=
	\begin{pmatrix}
		\psi_1 \\
		\psi_2 \\
		\psi_3 \\
		\psi_4
	\end{pmatrix}.
	\label{eq:dirac_spinor_spinstruktur}
\end{equation}

Die Komponenten sind durch die Dirac-Gleichung miteinander gekoppelt.
Sie enthalten unter anderem Informationen über die beiden möglichen
Spinrichtungen.

Im nichtrelativistischen Grenzfall werden die positiven
Energielösungen im Wesentlichen durch zwei Komponenten beschrieben.
Diese beiden Komponenten entsprechen den beiden möglichen
Spinorientierungen.

Vereinfacht kann man sie als

\begin{equation}
	\chi_{\uparrow}
	=
	\begin{pmatrix}
		1 \\
		0
	\end{pmatrix}
\end{equation}

und

\begin{equation}
	\chi_{\downarrow}
	=
	\begin{pmatrix}
		0 \\
		1
	\end{pmatrix}
\end{equation}

darstellen.

Diese zweikomponentigen Größen werden als Pauli-Spinoren%
\index{Pauli-Spinor} bezeichnet.

\begin{PhysicsBox}[Zwei mögliche Spinrichtungen]
	\label{box:zwei_spinrichtungen_dirac}
	
	Für eine Messung des Spins entlang einer festgelegten Raumrichtung
	gibt es beim Elektron zwei mögliche Ergebnisse:
	
	\begin{equation}
		s_z
		=
		+\frac{\hbar}{2}
	\end{equation}
	
	oder
	
	\begin{equation}
		s_z
		=
		-\frac{\hbar}{2}.
	\end{equation}
	
	Diese beiden Möglichkeiten werden häufig als Spin nach oben und Spin
	nach unten bezeichnet.
	
	Sie sind keine klassischen räumlichen Ausrichtungen einer rotierenden
	Kugel, sondern mögliche Messergebnisse eines Quantenzustandes.
\end{PhysicsBox}

\subsubsection*{Der Spinoperator}

In der Quantenmechanik werden messbare Größen durch Operatoren
beschrieben.

Der Spinoperator des Elektrons kann in der Dirac-Theorie in der Form

\begin{equation}
	\hat{\mathbf{S}}
	=
	\frac{\hbar}{2}
	\boldsymbol{\Sigma}
	\label{eq:spinoperator_dirac}
\end{equation}

geschrieben werden.

Die Größe \(\boldsymbol{\Sigma}\) besteht aus Matrizen, die auf die
Komponenten des Dirac-Spinors wirken.

In einer gebräuchlichen Darstellung gilt

\begin{equation}
	\Sigma_i
	=
	\begin{pmatrix}
		\sigma_i & 0 \\
		0 & \sigma_i
	\end{pmatrix},
	\label{eq:sigma_dirac_spinoperator}
\end{equation}

wobei \(\sigma_i\) die Pauli-Matrizen%
\index{Pauli-Matrizen!Spinoperator} sind.

Für das grundlegende Verständnis ist vor allem entscheidend, dass der
Spinoperator unmittelbar aus derselben Matrizenstruktur entsteht, die
bereits für die Dirac-Gleichung notwendig war.

\begin{MathBox}[Spinoperator der Dirac-Theorie]
	\label{box:spinoperator_dirac}
	
	Der Spinoperator lautet
	
	\begin{equation}
		\hat{\mathbf{S}}
		=
		\frac{\hbar}{2}
		\boldsymbol{\Sigma}.
	\end{equation}
	
	Der Faktor \(\hbar/2\) zeigt, dass das Elektron ein Teilchen mit
	Spin \(1/2\) ist.
	
	Die Matrizen \(\boldsymbol{\Sigma}\) wirken auf die Komponenten des
	Dirac-Spinors und unterscheiden die möglichen Spinrichtungen.
\end{MathBox}

\subsubsection*{Spin und Gesamtdrehimpuls}

In einem System können Bahndrehimpuls und Spin gleichzeitig auftreten.

Der gesamte Drehimpuls%
\index{Gesamtdrehimpuls!Elektron} setzt sich dann aus beiden Anteilen
zusammen:

\begin{equation}
	\mathbf{J}
	=
	\mathbf{L}
	+
	\mathbf{S}.
	\label{eq:gesamtdrehimpuls_dirac_spin}
\end{equation}

Dabei bezeichnet

\begin{itemize}
	\item \(\mathbf{L}\) den Bahndrehimpuls und
	\item \(\mathbf{S}\) den Spin.
\end{itemize}

In der relativistischen Theorie ist der Bahndrehimpuls allein im
Allgemeinen keine erhaltene Größe. Erst der Gesamtdrehimpuls aus
Bahndrehimpuls und Spin besitzt die richtige Erhaltungseigenschaft.

Dies zeigt, dass beide Größen in der Dirac-Theorie eng miteinander
verbunden sind.

\begin{PhysicsBox}[Erhaltung des Gesamtdrehimpulses]
	\label{box:erhaltung_gesamtdrehimpuls_dirac}
	
	Für ein relativistisches Elektron müssen Bahndrehimpuls und Spin
	gemeinsam betrachtet werden:
	
	\begin{equation}
		\mathbf{J}
		=
		\mathbf{L}
		+
		\mathbf{S}.
	\end{equation}
	
	Die Dirac-Gleichung zeigt damit, dass der Spin ein notwendiger
	Bestandteil des gesamten Drehimpulses ist.
\end{PhysicsBox}

\subsubsection*{Der Spin im Magnetfeld}

Das Elektron besitzt eine elektrische Ladung. Sein Spin ist deshalb mit
einem magnetischen Moment%
\index{magnetisches Moment!Spin des Elektrons} verbunden.

In einem äußeren Magnetfeld \(\mathbf{B}\) besitzt dieses magnetische
Moment eine Wechselwirkungsenergie:

\begin{equation}
	E_{\mathrm{mag}}
	=
	-
	\boldsymbol{\mu}_{\mathrm s}
	\cdot
	\mathbf{B}.
	\label{eq:spin_wechselwirkungsenergie_magnetfeld}
\end{equation}

Je nach Spinrichtung relativ zum Magnetfeld entstehen unterschiedliche
Energien.

Dadurch können zuvor gleiche Energiezustände aufgespalten werden. Diese
Aufspaltung bildet die Grundlage zahlreicher spektroskopischer
Messverfahren.

Das magnetische Moment des Elektrons lautet

\begin{equation}
	\boldsymbol{\mu}_{\mathrm s}
	=
	-g
	\frac{e}{2m_{\mathrm e}}
	\mathbf{S}.
	\label{eq:magnetisches_moment_spin_dirac}
\end{equation}

Das Minuszeichen hängt mit der negativen elektrischen Ladung des
Elektrons zusammen.

\begin{DidacticBox}[Spin im Magnetfeld]
	\label{box:spin_im_magnetfeld}
	
	In einem Magnetfeld hängt die Energie eines Elektrons von seiner
	Spinrichtung ab.
	
	Die beiden möglichen Spinrichtungen können deshalb unterschiedliche
	Energien besitzen.
	
	Diese Energieaufspaltung kann experimentell gemessen werden und macht
	den Spin sichtbar.
\end{DidacticBox}

\subsubsection*{Der \(g\)-Faktor des Elektrons}

Der dimensionslose Faktor \(g\) bestimmt die Stärke der Kopplung
zwischen Spin und magnetischem Moment.

Aus der Dirac-Gleichung folgt

\begin{equation}
	g
	=
	2.
	\label{eq:g_faktor_aus_dirac_spin}
\end{equation}

Dieses Ergebnis war bemerkenswert, weil es nicht zusätzlich in die
Theorie eingesetzt werden musste.

Der experimentell gemessene Wert liegt geringfügig über \(2\). Die
kleine Abweichung wird als anomales magnetisches Moment%
\index{anomales magnetisches Moment!Elektron} bezeichnet.

Sie entsteht durch Quanteneffekte, die erst in der
Quantenelektrodynamik%
\index{Quantenelektrodynamik!anomales magnetisches Moment}
vollständig beschrieben werden.

\begin{PhysicsBox}[Diracs Vorhersage des \(g\)-Faktors]
	\label{box:dirac_vorhersage_g_faktor}
	
	Die Dirac-Gleichung liefert für ein punktförmiges Elektron
	
	\begin{equation}
		g=2.
	\end{equation}
	
	Dieser Wert folgt unmittelbar aus der relativistischen
	Spinorstruktur.
	
	Die geringe Abweichung des gemessenen Wertes von \(2\) ist keine
	Widerlegung der Dirac-Gleichung. Sie entsteht durch zusätzliche
	Quanteneffekte des elektromagnetischen Feldes.
\end{PhysicsBox}

\subsubsection*{Der nichtrelativistische Grenzfall}

Für Elektronen mit kleinen Geschwindigkeiten kann die Dirac-Gleichung
näherungsweise vereinfacht werden.

Dabei entsteht die Pauli-Gleichung%
\index{Pauli-Gleichung!aus der Dirac-Gleichung}. Sie enthält bereits
die Wechselwirkung zwischen Spin und Magnetfeld.

Vereinfacht lautet der zugehörige Hamiltonoperator

\begin{equation}
	\hat{H}_{\mathrm P}
	=
	\frac{1}{2m_{\mathrm e}}
	\left(
	\hat{\mathbf{p}}
	+
	e\mathbf{A}
	\right)^2
	-
	\boldsymbol{\mu}_{\mathrm s}
	\cdot
	\mathbf{B},
	\label{eq:pauli_hamiltonoperator_aus_dirac}
\end{equation}

wobei \(\mathbf{A}\) das magnetische Vektorpotential%
\index{Vektorpotential!magnetisches} bezeichnet.

Der zweite Term enthält unmittelbar die Wechselwirkung zwischen Spin
und Magnetfeld.

Damit zeigt sich:

\begin{equation}
	\text{Dirac-Gleichung}
	\longrightarrow
	\text{Pauli-Gleichung}
\end{equation}

für kleine Geschwindigkeiten.

Wird zusätzlich der Spin vernachlässigt, erhält man schließlich die
gewöhnliche Schrödingergleichung.

\begin{DidacticBox}[Die Pauli-Gleichung als Näherung]
	\label{box:pauli_gleichung_naeherung_dirac}
	
	Die Pauli-Gleichung ist keine unabhängige Erklärung des Spins.
	
	Sie ist die nichtrelativistische Näherung der Dirac-Gleichung für ein
	Teilchen mit Spin \(1/2\).
	
	Damit wird verständlich, warum der Spin in der Pauli-Gleichung
	zusätzlich erscheint, in der Dirac-Gleichung jedoch bereits enthalten
	ist.
\end{DidacticBox}

\subsubsection*{Warum der Spin zwangsläufig erscheint}

Dirac hatte den Spin nicht als Ausgangspunkt seiner Theorie verwendet.

Seine Forderungen waren zunächst rein physikalischer und mathematischer
Natur:

\begin{itemize}
	\item Die Gleichung sollte relativistisch sein.
	\item Sie sollte nur Ableitungen erster Ordnung enthalten.
	\item Sie sollte eine sinnvolle Wahrscheinlichkeitsdichte besitzen.
	\item Beim Quadrieren sollte die relativistische
	Energie-Impuls-Beziehung entstehen.
\end{itemize}

Diese Forderungen machten Matrizen und eine mehrkomponentige
Wellenfunktion notwendig.

Aus genau dieser Struktur entstand der Spin.

\begin{HistoryBox}[Der Spin als theoretische Notwendigkeit]
	\label{box:spin_theoretische_notwendigkeit}
	
	Dirac setzte den Elektronenspin nicht von Hand in seine Gleichung ein.
	
	Der Spin erschien, weil eine relativistische Quantengleichung für das
	Elektron eine mehrkomponentige Wellenfunktion und eine besondere
	Matrizenstruktur benötigt.
	
	Damit wurde aus einer zuvor zusätzlich angenommenen Eigenschaft eine
	Folge der Theorie.
\end{HistoryBox}

Die Dirac-Gleichung erklärte somit nicht nur die relativistische
Bewegung des Elektrons. Sie zeigte auch, warum ein Elektron zwei
Spinrichtungen besitzt und weshalb sein Spin mit einem magnetischen
Moment verbunden ist.

Neben dem Spin enthielt die Gleichung jedoch noch eine weitere
unerwartete Eigenschaft: Sie besaß Lösungen mit negativer Energie.
Deren Deutung führte schließlich zur Vorhersage des Positrons.

